\begin{abstract}
Micro, Small and Medium Enterprises are important contributors to industrial development, but many MSMEs face high energy consumption, high electricity cost, diesel-generator dependency and limited access to intelligent energy management systems. Traditional energy monitoring systems mainly display past or present readings and do not provide prediction, optimization or decision support. Because of this, industries may face energy wastage, peak-load penalties and unnecessary carbon dioxide emissions.

This project presents an AI-driven digital twin for smart energy optimization in MSMEs using IoT-based energy analytics. The system uses historical industrial energy data consisting of production output, grid energy and diesel-generator energy. The data is cleaned and processed using MATLAB, and important Key Performance Indicators such as total energy consumption, Energy Per Unit, energy efficiency, operating cost, generator dependency and carbon dioxide emissions are calculated.

Linear and quadratic regression models are developed to study the relationship between production and total energy consumption. The quadratic and linear regression models are selected as the optimal modeling approach because they provide stable coefficients of determination without overfitting the limited dataset. While complex machine learning techniques such as Random Forest, XGBoost and LSTM forecasting were evaluated, they were excluded from the final digital twin due to severe overfitting and test-leakage caused by the restricted sample size (59 records).

A digital twin simulation engine is developed to test different operational scenarios without disturbing the real manufacturing system. The scenarios include mathematical energy baseline bounding, diesel-generator reduction, production scaling, base-load reduction and peak-load avoidance. A dashboard-based decision support system is proposed to display KPIs, energy trends, anomaly alerts, load scheduling recommendations and what-if simulation results.

The results show that the AI-driven digital twin can support MSMEs in reducing energy cost, improving energy efficiency, decreasing diesel-generator usage and supporting sustainable industrial operation. The project demonstrates how Industry 4.0 technologies such as IoT, AI and digital twins can be used as a low-cost and scalable solution for smart energy management in MSMEs.
\end{abstract}
